



\section{Theorems about transifinite- and epsilon- induction and recursion.}

\begin{recall}
Transfinite recursion\\
$G:V\to V$\\
$\exists! F:\Ord\to V$, s.t.\\
$F(a) = G(F\restriction a)$\\
\end{recall}
%Recall:\\
\begin{recall}
Epsilon recursion:\\
$G:V\to V$\\
$\exists! F:V\to V$, s.t.\\
$F(a) = G(F\restriction a)$\\
\end{recall}

\begin{definition}
Height function definition:\\
$h(x)$ is the height of a the set. It is equal by definition to the maximum length of a $\ni$-chain, started from $x$.\\
\end{definition}

Examples:\\
$h(\emptyset)$=0,\\
$h(\{\emptyset\})$=1,\\
$h(\{\emptyset,\{\emptyset\}\})$=2,\\
and in general if $n$ is a finite ordinal number, then $h(n)=n$.\\
Definition: by $\in$-recursion
\[
h(x) := \sup \{ h(y) | y \in x\}
\]

Property:
\[
h(x) = \sup h[x]
\]

$G(R) := \sup(\ran(R)) = \union(\ran(R))$\\
via epsilon recursion.\\
$h : V\to V$\\
$h(a) = \sup(\ran(h\restriction a) )$\\



\begin{lemma}()
\[
(h:V\to\Ord)
\]
\end{lemma}
\Proof
by epsilon induction.\\
\Done

\subsection{older Inductions and recursions thms}%
There exist unique class-function rec, such that:
\[
\mbox{rec}:(V\times V\to V)\to(V\to V)
\]
\[
\mbox{rec}(H) = H\ \circ\ (\mbox{rec}(H)\restriction_-)
\]
\pagebreak

